% Appendix: QA records
\chapter*{附录：问答记录}
\label{chap:Appendix}

本节记录学习 Do Carmo《曲线与曲面的微分几何》过程中的问答与思考。

% 待填充内容
\subsection*{待记录问题}

本节将随学习进度逐步填充以下问题：

\begin{itemize}
\item 曲率与挠率的直观几何意义
\item Frenet-Serret 公式的推导细节
\item 局部标准形式的系数解释
\item 全曲率定理的证明思路
\item 等周不等式的证明
\end{itemize}

\subsection*{历史悠久的问答}

% 从以下问题开始记录：

\subsection*{曲率的直观理解}\label{sec:qa-curvature-intuition}

\textbf{问}：曲率 $k(s) = \|\alpha''(s)\|$ 的直观几何意义是什么？

\textbf{答}：对于单位速度曲线 $\alpha(s)$，速度向量 $\alpha'(s)$ 是单位切向量，其模长恒为1。加速度向量 $\alpha''(s)$ 描述速度向量方向的变化率。 $\|\alpha''(s)\|$ 正是这种方向变化的速率——即曲线在每一点的弯曲程度。直线速度方向不变，$\alpha''(s) = 0$；圆周速度方向匀速变化，$\|\alpha''(s)\| = 1/a$（$a$ 为半径）。

\subsection*{挠率与扭曲}\label{sec:qa-torsion}

\textbf{问}：挠率 $\tau(s)$ 描述的是什么？

\textbf{答}：挠率描述曲线在空间中的"扭曲"程度。曲率只告诉我们曲线在密切平面内如何弯曲，但曲线可能"偏离"密切平面——这种偏离的速度由挠率度量。对于平面曲线，$\alpha', \alpha'', \alpha'''$ 始终共面，$\det(\alpha', \alpha'', \alpha''') = 0$，故挠率为零。螺旋线的挠率为常数，描述其均匀扭曲的特性。

\subsection*{曲率另一个表达式的推导}\label{sec:appendix-ch1-curvature-alt}

\textbf{背景}：命题 1-5, 2 给出了曲率在一般参数化下的表达式：
\[
k(t) = \frac{\|\alpha'(t) \times \alpha''(t)\|}{\|\alpha'(t)\|^3}.
\]

\textbf{目标}：从单位速度曲率的定义 $k(s) = \|\alpha''(s)\|$ 推导出一般参数下的表达式。

\textbf{推导步骤}：

设 $\alpha(t)$ 是正则曲线，$s = s(t)$ 是弧长参数。回忆弧长公式：
\[
s(t) = \int_{t_0}^t \|\alpha'(u)\| \, du, \quad \frac{ds}{dt} = \|\alpha'(t)\|.
\]

由链式法则，$\alpha'(s) = \dfrac{\alpha'(t)}{ds/dt} = \dfrac{\alpha'(t)}{\|\alpha'(t)\|}$。

进一步，
\[
\alpha''(s) = \frac{d}{ds}\alpha'(s) = \frac{d}{dt}\left(\frac{\alpha'(t)}{\|\alpha'(t)\|}\right) \cdot \frac{dt}{ds}.
\]

经过计算（利用向量积的模长公式 $\|u \times v\| = \|u\|\|v\|\sin\theta$），可得：
\[
\|\alpha''(s)\| = \frac{\|\alpha'(t) \times \alpha''(t)\|}{\|\alpha'(t)\|^3}.
\]

因此在参数 $t$ 下，曲率为 $k(t) = \dfrac{\|\alpha'(t) \times \alpha''(t)\|}{\|\alpha'(t)\|^3}$。

\subsection*{挠率另一个表达式的推导}\label{sec:appendix-ch1-torsion-alt}

\textbf{背景}：命题 1-5, 4 给出了挠率的另一个表达式：
\[
\tau(s) = -\langle \mathbf{b}'(s), \mathbf{n}(s) \rangle.
\]

\textbf{目标}：从挠率的行列式定义 $\tau = \dfrac{\det(\alpha', \alpha'', \alpha''')}{k^2}$ 推导出 $\tau = -\langle \mathbf{b}', \mathbf{n} \rangle$。

\textbf{推导步骤}：

由 Frenet 标架的定义：$\mathbf{t} = \alpha'$，$\mathbf{n} = \dfrac{\alpha''}{k}$，$\mathbf{b} = \mathbf{t} \times \mathbf{n}$。

对 $\mathbf{b} = \mathbf{t} \times \mathbf{n}$ 求导：
\[
\mathbf{b}' = \mathbf{t}' \times \mathbf{n} + \mathbf{t} \times \mathbf{n}' = (k\mathbf{n}) \times \mathbf{n} + \mathbf{t} \times \mathbf{n}' = \mathbf{t} \times \mathbf{n}'.
\]

由于 $\mathbf{n}'$ 可以分解为 $A\mathbf{t} + B\mathbf{n} + C\mathbf{b}$（由正交性），且 $\|\mathbf{n}\|=1$ 蕴含 $\langle \mathbf{n}', \mathbf{n} \rangle = 0$，故 $\mathbf{n}' = -k\mathbf{t} + \tau\mathbf{b}$（这正是 Frenet-Serret 公式的第二式）。

于是
\[
\mathbf{b}' = \mathbf{t} \times (-k\mathbf{t} + \tau\mathbf{b}) = \mathbf{t} \times \tau\mathbf{b} = \tau(\mathbf{t} \times \mathbf{b}) = -\tau\mathbf{n},
\]
因为 $\mathbf{t} \times \mathbf{b} = -\mathbf{b} \times \mathbf{t} = -\mathbf{n}$。

因此 $\mathbf{b}' = -\tau\mathbf{n}$，取与 $\mathbf{n}$ 的内积即得
\[
\tau = -\langle \mathbf{b}', \mathbf{n} \rangle.
\]

\subsection*{Frenet-Serret 公式的直观解释与推导思路}\label{sec:appendix-ch1-frenet-serret}

\textbf{背景}：Frenet-Serret 公式是曲线局部理论的核心：
\[
\begin{cases}
\mathbf{t}' = k \mathbf{n}, \\
\mathbf{n}' = -k \mathbf{t} + \tau \mathbf{b}, \\
\mathbf{b}' = -\tau \mathbf{n}.
\end{cases}
\]

\textbf{直观解释}：

可以将 Frenet 标架 $\{\mathbf{t}, \mathbf{n}, \mathbf{b}\}$ 想象为沿曲线运动的"观察者"的坐标系。这个坐标系随着曲线在空间中移动而转动和摆动。

\begin{itemize}
\item $\mathbf{t}' = k\mathbf{n}$：切向量的变化完全由法向分量决定，系数 $k$ 正是弯曲程度——$k$ 越大，切向量转动得越快。

\item $\mathbf{n}' = -k\mathbf{t} + \tau\mathbf{b}$：主法向量的变化涉及两个方面：$-k\mathbf{t}$ 部分反映了主法向量试图"追回"切向量方向变化（与 $\mathbf{t}' = k\mathbf{n}$ 相呼应），而 $\tau\mathbf{b}$ 部分则反映了曲线在空间中的扭曲。

\item $\mathbf{b}' = -\tau\mathbf{n}$：副法向量的变化只与挠率有关——它描述了密切平面的旋转。
\end{itemize}

\textbf{几何图像}：想象一个骑自行车的人沿着曲线运动。人的朝向就是 $\mathbf{t}$ 方向，车把的指向是 $\mathbf{n}$，而 $\mathbf{b}$ 则垂直于前两者。曲率 $k$ 决定了人在转弯时需要倾斜多少（这对应 $\mathbf{t}' = k\mathbf{n}$），而挠率 $\tau$ 则描述了车子在转弯的同时是否还在上下颠簸（对应 $\mathbf{b}' = -\tau\mathbf{n}$）。

\textbf{推导思路}（简述）：

Frenet-Serret 公式可从正交性条件出发推导。由于 $\{\mathbf{t}, \mathbf{n}, \mathbf{b}\}$ 是标准正交基，有 $\langle \mathbf{t}, \mathbf{n} \rangle = 0$，$\langle \mathbf{t}, \mathbf{b} \rangle = 0$，$\langle \mathbf{n}, \mathbf{b} \rangle = 0$，以及 $\|\mathbf{t}\| = \|\mathbf{n}\| = \|\mathbf{b}\| = 1$。

对正交性条件求导，并利用 $\mathbf{b} = \mathbf{t} \times \mathbf{n}$，可以逐步得到三个公式。例如：

\begin{enumerate}
\item 由 $\langle \mathbf{t}, \mathbf{n} \rangle = 0$ 求导：$\langle \mathbf{t}', \mathbf{n} \rangle + \langle \mathbf{t}, \mathbf{n}' \rangle = 0$。
\item 由于 $\mathbf{t}' = \alpha''$ 与 $\mathbf{n}$ 平行（定义 $\mathbf{n} = \mathbf{t}'/|\mathbf{t}'|$），得 $\langle \mathbf{t}', \mathbf{n} \rangle = |\mathbf{t}'| = k$。
\item 结合 $\|\mathbf{n}'\|$ 的计算，可得 $\mathbf{t}' = k\mathbf{n}$ 和 $\mathbf{n}' = -k\mathbf{t} + \tau\mathbf{b}$。
\end{enumerate}

详细证明见 do Carmo, Proposition 1-5, 4。
